% 05-lexical.tex — Clause 5: Lexical conventions
\chapter{Lexical conventions}
\label{sec:05-lexical}
\index{lexical conventions}

\pnum
A Lain source file is translated in two conceptual steps: (1) the source
text is decomposed into tokens as described in this clause; (2) the token
stream is parsed into an abstract syntax tree according to the grammar.
Implementation may combine these steps.

\section{Source file encoding}
\label{sec:lex.encoding}
\index{source file!encoding}

\pnum
A Lain source file shall be a sequence of bytes encoded in UTF-8 (Unicode
Standard, Annex D).  The byte order mark (U+FEFF) shall not appear in a
source file.  If it does, behavior is \idbehav{whether a leading BOM is
silently skipped or causes a diagnostic}.

\pnum
Multi-byte UTF-8 sequences are permitted inside string literals and comments
but shall not appear in identifiers.  Identifiers shall consist solely of
ASCII characters (see \S\,\ref{sec:lex.identifiers}).

\pnum
Line endings are normalized: a carriage return (U+000D) followed
immediately by a line feed (U+000A) is treated as a single line feed.
A lone carriage return is treated as a line feed.

\section{Comments}
\label{sec:lex.comments}
\index{comment}

\pnum
A \textit{line comment} begins with \tok{//} and extends to the first
following line feed character.  The line feed is not part of the comment.

\pnum
A \textit{block comment} begins with \tok{/*} and ends with the nearest
following \tok{*/}.  Block comments may be nested: each opening \tok{/*}
inside a block comment opens a new nesting level requiring its own closing
\tok{*/}.

\pnum
Comments are stripped during lexical analysis and have no effect on the
semantics of the enclosing program.

\begin{constraint}
  An unterminated block comment (one for which no matching \tok{*/} is
  found before the end of the source file) is \illformed{}~(\diagcode{E100}).
\end{constraint}

\section{Tokens}
\label{sec:lex.tokens}
\index{token}

\pnum
The result of lexical analysis is a sequence of tokens.  Whitespace and
comments serve solely to separate tokens.  The grammar production for
\gramref{token} is given in Annex~A (\S\,\ref{gram:token}).

\pnum
Where the token stream is ambiguous, the lexer shall apply the
\textit{maximal munch} rule: at each position, the longest sequence of
characters that forms a valid token is consumed.

\begin{example}
\begin{laincode}
..=    // one token: range-inclusive, NOT '.' followed by '.='
>=     // one token: greater-than-or-equal, NOT '>' followed by '='
\end{laincode}
\end{example}

\pnum
Keywords (\S\,\ref{sec:lex.keywords}) take priority over identifiers: the
character sequence \texttt{func} at a point where an identifier is expected
is always the keyword \tok{func}.

\section{Whitespace}
\label{sec:lex.whitespace}
\index{whitespace}

\pnum
The whitespace characters are: space (U+0020) and horizontal tab (U+0009).
They serve only to separate adjacent tokens.

\pnum
Line feeds (U+000A), after normalization (\S\,\ref{sec:lex.encoding}), are
significant: they act as statement terminators in contexts where a statement
may follow (see \S\,\ref{sec:09-statements}).  Multiple consecutive line
feeds are treated as a single statement boundary.

\begin{constraint}
  A semicolon (\tok{;}) shall not appear in a Lain source file.  If one is
  present, the implementation shall emit~\diagcode{E100}.
\end{constraint}

\begin{rationale}
  Newlines as statement terminators eliminate the class of bugs caused by
  missing or accidental semicolons.  They also prevent multi-statement lines,
  keeping statement boundaries unambiguous.
\end{rationale}

\section{Keywords}
\label{sec:lex.keywords}
\index{keyword}

\pnum
The following identifiers are reserved as keywords.  They shall not be
used as user-defined identifiers.  The grammar production is
\gramref{keyword} (Annex~A).

\begin{longtable}{@{}lll@{}}
\toprule
\textbf{Keyword} & \textbf{Purpose} & \textbf{Clause} \\
\midrule
\endfirsthead
\multicolumn{3}{l}{\textit{(continued)}}\\
\toprule
\textbf{Keyword} & \textbf{Purpose} & \textbf{Clause} \\
\midrule
\endhead
\bottomrule
\endlastfoot
\tok{and}         & Logical AND operator           & \S\,\ref{sec:08-expressions} \\
\tok{as}          & Type cast                      & \S\,\ref{sec:08-expressions} \\
\tok{case}        & Pattern matching               & \S\,\ref{sec:15-match} \\
\tok{comptime}    & Compile-time evaluation        & \S\,\ref{sec:14-generics} \\
\tok{decreasing}  & Bounded-loop termination measure & \S\,\ref{sec:09-statements} \\
\tok{defer}       & Deferred execution             & \S\,\ref{sec:09-statements} \\
\tok{else}        & Alternative branch             & \S\,\ref{sec:09-statements} \\
\tok{false}       & Boolean literal                & \S\,\ref{sec:lex.literals} \\
\tok{for}         & Range-based iteration          & \S\,\ref{sec:09-statements} \\
\tok{func}        & Pure function declaration      & \S\,\ref{sec:12-functions} \\
\tok{if}          & Conditional                    & \S\,\ref{sec:09-statements} \\
\tok{import}      & Module import                  & \S\,\ref{sec:16-modules} \\
\tok{in}          & Range iteration / bounds guard & \S\,\ref{sec:08-expressions} \\
\tok{module}      & Module declaration             & \S\,\ref{sec:16-modules} \\
\tok{mov}         & Ownership transfer             & \S\,\ref{sec:11-ownership} \\
\tok{mut}         & Mutable borrow expression      & \S\,\ref{sec:11-ownership} \\
\tok{not}         & Logical NOT                    & \S\,\ref{sec:08-expressions} \\
\tok{or}          & Logical OR operator            & \S\,\ref{sec:08-expressions} \\
\tok{proc}        & Procedure declaration          & \S\,\ref{sec:12-functions} \\
\tok{ref}         & Shared borrow expression       & \S\,\ref{sec:11-ownership} \\
\tok{return}      & Return statement               & \S\,\ref{sec:09-statements} \\
\tok{struct}      & Struct type declaration        & \S\,\ref{sec:10-declarations} \\
\tok{enum}        & Enum type declaration          & \S\,\ref{sec:10-declarations} \\
\tok{true}        & Boolean literal                & \S\,\ref{sec:lex.literals} \\
\tok{unsafe}      & Unsafe block                   & \S\,\ref{sec:18-unsafe} \\
\tok{var}         & Mutable binding / mutable borrow & \S\,\ref{sec:10-declarations} \\
\tok{while}       & Loop statement                 & \S\,\ref{sec:09-statements} \\
\end{longtable}

\pnum
The following identifiers are \textit{reserved keywords}: they are
recognized by the lexer but have no current semantics.  They are reserved
for future use.

\begin{longtable}{@{}ll@{}}
\toprule
\textbf{Reserved keyword} & \textbf{Intended future purpose} \\
\midrule
\endfirsthead
\multicolumn{2}{l}{\textit{(continued)}}\\
\toprule
\textbf{Reserved keyword} & \textbf{Intended future purpose} \\
\midrule
\endhead
\bottomrule
\endlastfoot
\tok{end}     & Block terminator \\
\tok{export}  & Module visibility control \\
\tok{extern}  & C interoperability declarations \\
\tok{type}    & Type alias / type parameter \\
\end{longtable}

\begin{constraint}
  An identifier that is identical to a keyword or reserved keyword shall
  not be used in any context that requires an identifier.  Such use is
  \illformed{}~(\diagcode{E100}).
\end{constraint}

\section{Identifiers}
\label{sec:lex.identifiers}
\index{identifier}

\pnum
An identifier consists of ASCII letters, decimal digits, and underscores,
and shall begin with a letter or underscore.  The grammar production is
\gramref{identifier} (Annex~A).

\pnum
Identifiers are case-sensitive: \texttt{foo}, \texttt{Foo}, and
\texttt{FOO} are three distinct identifiers.

\pnum
A conforming implementation shall support identifiers of at least 255
characters.  The behavior when an identifier exceeds this length is
\idbehav{whether oversized identifiers are truncated or cause a diagnostic}.

\begin{constraint}
  An identifier shall not be identical to a keyword
  (\S\,\ref{sec:lex.keywords}).
\end{constraint}

\begin{constraint}
  An identifier shall begin with a letter or underscore, not a digit.
  A sequence of characters that begins with a digit is an integer
  literal~(\S\,\ref{sec:lex.literals}), not an identifier.
\end{constraint}

\section{Literals}
\label{sec:lex.literals}
\index{literal}

\subsection{Integer literals}
\label{sec:lex.literals.integer}
\index{integer literal}

\pnum
An integer literal denotes a constant integer value.  Four notations are
supported (grammar: \gramref{integer-literal}):

\begin{longtable}{@{}lll@{}}
\toprule
\textbf{Notation} & \textbf{Prefix} & \textbf{Example} \\
\midrule
\endfirsthead
\toprule
\textbf{Notation} & \textbf{Prefix} & \textbf{Example} \\
\midrule
\endhead
\bottomrule
\endlastfoot
Decimal   & (none) & \lain{42}, \lain{1_000_000} \\
Hexadecimal & \lain{0x} & \lain{0xFF}, \lain{0x1A_3B} \\
Binary    & \lain{0b} & \lain{0b1010_0101} \\
Octal     & \lain{0o} & \lain{0o755} \\
\end{longtable}

\pnum
An underscore (\texttt{\_}) may appear between any two digits of an integer
literal as a visual separator.  It has no effect on the value.

\pnum
The type of an integer literal is inferred from context (see
\S\,\ref{sec:07-types}).  In the absence of type context, the default type
is \lain{int}.

\begin{constraint}
  An integer literal whose value does not fit in the type determined by
  context is \illformed{}~(\diagcode{E100}).
\end{constraint}

\begin{example}
\begin{laincode}
var a = 255         // int, value 255
var b u8 = 255      // u8, value 255 -- OK
var c u8 = 256      // [E100]: 256 does not fit in u8
var d = 0xFF        // int, value 255
var e = 0b1111_0000 // int, value 240
var f = 1_000_000   // int, value 1000000
\end{laincode}
\end{example}

\subsection{Floating-point literals}
\label{sec:lex.literals.float}
\index{floating-point literal}

\pnum
A floating-point literal consists of a decimal integer part, a decimal
point, and a fractional part.  An optional exponent part
\texttt{e}\textit{N} or \texttt{E}\textit{N} may follow (grammar:
\gramref{float-literal}).

\pnum
The default type of a floating-point literal is \lain{f64}.

\begin{example}
\begin{laincode}
var x = 3.14      // f64
var y f32 = 1.0   // f32
var z = 2.5e3     // f64, value 2500.0
\end{laincode}
\end{example}

\subsection{Boolean literals}
\label{sec:lex.literals.bool}
\index{boolean literal}

\pnum
The boolean literals are \tok{true} and \tok{false} (grammar:
\gramref{bool-literal}).  Both have type \lain{bool}.

\subsection{Character literals}
\label{sec:lex.literals.char}
\index{character literal}

\pnum
A character literal is a single byte value enclosed in single quotes.  It
has type \lain{u8}.  Escape sequences (\S\,\ref{sec:lex.escape}) are
permitted.

\begin{example}
\begin{laincode}
var a = 'A'    // u8, value 65
var n = '\n'   // u8, value 10
var z = '\0'   // u8, value 0
\end{laincode}
\end{example}

\begin{constraint}
  A character literal shall contain exactly one character or one escape
  sequence.  An empty or multi-character character literal is
  \illformed{}~(\diagcode{E100}).
\end{constraint}

\subsection{String literals}
\label{sec:lex.literals.string}
\index{string literal}

\pnum
A string literal is a sequence of zero or more characters enclosed in
double quotes.  Escape sequences (\S\,\ref{sec:lex.escape}) are permitted
(grammar: \gramref{string-literal}).

\pnum
The type of a string literal of $N$ visible characters is
\lain{Fixed_u8_}\textit{N}\texttt{+1} (an array of $N{+}1$ bytes,
null-terminated).  This type is implicitly coercible to \lain{[u8]} (a
byte slice).

\pnum
The \texttt{.len} property of a string literal value equals $N$ (excluding
the null terminator).

\begin{example}
\begin{laincode}
var s = "hello"   // type: Fixed_u8_6 (5 chars + null)
                  // s.len == 5
\end{laincode}
\end{example}

\pnum
String literals are stored in the read-only data segment of the translation
output.  Modifying a string literal through a pointer is \undefbehav{the
result of writing through a pointer derived from a string literal}.

\subsection{Escape sequences}
\label{sec:lex.escape}
\index{escape sequence}

\pnum
The following escape sequences are recognized in string and character
literals:

\begin{longtable}{@{}lll@{}}
\toprule
\textbf{Escape} & \textbf{Byte value} & \textbf{Description} \\
\midrule
\endfirsthead
\toprule
\textbf{Escape} & \textbf{Byte value} & \textbf{Description} \\
\midrule
\endhead
\bottomrule
\endlastfoot
\texttt{\textbackslash n}  & 0x0A & Line feed (newline) \\
\texttt{\textbackslash t}  & 0x09 & Horizontal tab \\
\texttt{\textbackslash r}  & 0x0D & Carriage return \\
\texttt{\textbackslash 0}  & 0x00 & Null byte \\
\texttt{\textbackslash\textbackslash} & 0x5C & Backslash \\
\texttt{\textbackslash "} & 0x22 & Double quote \\
\texttt{\textbackslash '} & 0x27 & Single quote \\
\end{longtable}

\begin{constraint}
  A backslash not followed by one of the characters in the table above is
  an unrecognized escape sequence and is \illformed{}~(\diagcode{E100}).
\end{constraint}

\section{Operators and punctuators}
\label{sec:lex.operators}
\index{operator}

\pnum
The complete set of operator and punctuator tokens is listed below.  The
grammar productions \gramref{operator} and \gramref{punctuator} are in
Annex~A.

\subsection{Arithmetic operators}

\begin{longtable}{@{}ll@{}}
\toprule
\textbf{Token} & \textbf{Meaning} \\
\midrule
\endfirsthead
\toprule
\textbf{Token} & \textbf{Meaning} \\
\midrule
\endhead
\bottomrule
\endlastfoot
\tok{+}  & Addition; unary identity \\
\tok{-}  & Subtraction; unary negation \\
\tok{*}  & Multiplication \\
\tok{/}  & Division \\
\tok{\%} & Modulo (remainder) \\
\end{longtable}

\subsection{Comparison operators}

\begin{longtable}{@{}ll@{}}
\toprule
\textbf{Token} & \textbf{Meaning} \\
\midrule
\endfirsthead
\toprule
\textbf{Token} & \textbf{Meaning} \\
\midrule
\endhead
\bottomrule
\endlastfoot
\tok{==} & Equal \\
\tok{!=} & Not equal \\
\tok{<}  & Less than \\
\tok{>}  & Greater than \\
\tok{<=} & Less than or equal \\
\tok{>=} & Greater than or equal \\
\end{longtable}

\subsection{Logical operators}

\begin{longtable}{@{}ll@{}}
\toprule
\textbf{Token} & \textbf{Meaning} \\
\midrule
\endfirsthead
\toprule
\textbf{Token} & \textbf{Meaning} \\
\midrule
\endhead
\bottomrule
\endlastfoot
\tok{and} & Logical AND (short-circuit) \\
\tok{or}  & Logical OR (short-circuit) \\
\tok{not} & Logical NOT (prefix) \\
\tok{!}   & Logical NOT (alternate prefix form) \\
\end{longtable}

\subsection{Bitwise operators}

\begin{longtable}{@{}ll@{}}
\toprule
\textbf{Token} & \textbf{Meaning} \\
\midrule
\endfirsthead
\toprule
\textbf{Token} & \textbf{Meaning} \\
\midrule
\endhead
\bottomrule
\endlastfoot
\tok{\&}   & Bitwise AND \\
\tok{|}    & Bitwise OR \\
\tok{\^{}} & Bitwise XOR \\
\tok{<<}   & Left shift \\
\tok{>>}   & Right shift \\
\end{longtable}

\subsection{Assignment operators}

\begin{longtable}{@{}ll@{}}
\toprule
\textbf{Token} & \textbf{Meaning} \\
\midrule
\endfirsthead
\toprule
\textbf{Token} & \textbf{Meaning} \\
\midrule
\endhead
\bottomrule
\endlastfoot
\tok{=}    & Assignment \\
\tok{+=}   & Add and assign \\
\tok{-=}   & Subtract and assign \\
\tok{*=}   & Multiply and assign \\
\tok{/=}   & Divide and assign \\
\tok{\%=}  & Modulo and assign \\
\tok{\&=}  & Bitwise AND and assign \\
\tok{|=}   & Bitwise OR and assign \\
\tok{\^{}=}& Bitwise XOR and assign \\
\tok{<<=}  & Left shift and assign \\
\tok{>>=}  & Right shift and assign \\
\end{longtable}

\subsection{Other operators and punctuators}

\begin{longtable}{@{}ll@{}}
\toprule
\textbf{Token} & \textbf{Meaning} \\
\midrule
\endfirsthead
\toprule
\textbf{Token} & \textbf{Meaning} \\
\midrule
\endhead
\bottomrule
\endlastfoot
\tok{.}    & Member access; module path separator \\
\tok{..}   & Exclusive range (half-open interval) \\
\tok{->}   & Function return type separator \\
\tok{(}    & Open parenthesis \\
\tok{)}    & Close parenthesis \\
\tok{\{}   & Open brace \\
\tok{\}}   & Close brace \\
\tok{[}    & Open bracket \\
\tok{]}    & Close bracket \\
\tok{,}    & List separator \\
\tok{:}    & Type annotation; match arm separator \\
\tok{@}    & Attribute / annotation prefix \\
\end{longtable}

\section{Token disambiguation}
\label{sec:lex.disambiguation}
\index{token disambiguation}

\pnum
Several token sequences have context-dependent meaning that is resolved by
the parser, not the lexer.  The lexer produces the same token regardless of
context:

\begin{longtable}{@{}lll@{}}
\toprule
\textbf{Symbol} & \textbf{Context 1} & \textbf{Context 2} \\
\midrule
\endfirsthead
\toprule
\textbf{Symbol} & \textbf{Context 1} & \textbf{Context 2} \\
\midrule
\endhead
\bottomrule
\endlastfoot
\tok{*} & Multiplication (binary)            & (future: pointer dereference) \\
\tok{-} & Subtraction (binary)               & Negation (unary prefix) \\
\tok{.} & Member access                      & Module path separator \\
\tok{\&}& Borrowed match prefix              & Bitwise AND \\
\end{longtable}

\section{Integer arithmetic and overflow}
\label{sec:lex.overflow}
\index{integer overflow}

\pnum
Integer arithmetic in Lain uses two's complement representation for all
signed integer types.  Overflow wraps modulo $2^N$ where $N$ is the bit
width of the type.

\pnum
Unsigned integer arithmetic is modular (wrapping) by definition.

\pnum
A conforming implementation that emits C99 shall compile the translation
output with \texttt{-fwrapv} or equivalent, ensuring that signed integer
overflow is two's complement wrapping and not undefined behavior in the C
sense.

\begin{note}
  This differs from ISO C, where signed integer overflow is undefined
  behavior.  In Lain, signed overflow is defined as wrapping.  Programs
  that rely on wrapping behavior are well-formed.
\end{note}
